\documentclass[./2025_01.tex]{subfiles}

\begin{document}
\begin{center}
\makebox[\textwidth]{Nombre y Apellido:\enspace\hrulefill}
\end{center}
\begin{questions}
    \question Supóngase que se tiene una máquina térmica capaz de funcionar de los siguientes modos:¿Cuál es rendimiento del ciclo del ciclo ideal de Carnot que podría tener en cada caso?
    \begin{parts}
        \part[1] Entre \qty{200}{\degree\kelvin}  y \SI{400}{\degree\kelvin},
        \part[1] Entre \SI{600}{\degree\kelvin} y \SI{400}{\degree\kelvin}.
        \part[1] Entre \SI{100}{\celsius} y \SI{20}{\celsius}
        \part[1] Entre \SI{473}{\celsius} y \SI{683}{\celsius}
        
    \end{parts}

    \question Una máquina térmica con eficiencia de $\eta_{\%}=22\%$ produce $W=1530$J de trabajo. Encontrar: 
    \begin{parts}
        \part[2] La cantidad de calor absorbida del depósito térmico $Q_1$.
        \part[2] La cantidad de calor desechado al depósito térmico.$Q_2$
        \part[2] Las temperaturas del depósito de la fuente fría y fuente caliente $T_1$ y $T_2$.
    \end{parts} 
    
    
\end{questions}




 \end{document}